% note for Marginal_Effects_Panels.pdf (generated by R/41_Paper_Tables.R, PLLM-0256)
% caption (Title Case): \caption{Marginal Effect of Populism on Liberalism across the Ideology Scale, by Liberalism Dimension}
\par\smallskip\begin{minipage}{\linewidth}\scriptsize\textit{Notes:} Marginal effect of a one-point increase in populism on each liberalism measure (0--10), from column (5) of Tables 4--8 (populism, ideology and their interaction; country and year fixed effects; each outcome on its own sample, N between 2,999 and 3,192). The line traces the effect over the ideology score $(R-L)/(R+L)$ from $-$1 (left) to +1 (right), the shaded band is its confidence interval, the rug marks the observed ideology scores, the dashed red line is the populism coefficient of column (1) of the same table, and the text gives the $p$-value of the interaction. 95 percent confidence intervals from standard errors clustered by party (CMP party code), critical value of the $t$ distribution with the number of party clusters minus one degrees of freedom.\end{minipage}
